\documentclass[11pt]{article}
\PassOptionsToPackage{hyphens}{url}  % let long URLs break at hyphens/slashes

% --- Dependency-light preamble: compiles on a stock TeX Live / Overleaf / arXiv ---
\usepackage[utf8]{inputenc}
\usepackage[T1]{fontenc}
\usepackage{lmodern}
\usepackage[margin=1in]{geometry}
\usepackage{booktabs}
\usepackage{amsmath}
\usepackage{array}
\usepackage{microtype}
\usepackage[hidelinks]{hyperref}
\usepackage{url}
\usepackage{textcomp}
\setlength{\emergencystretch}{3em}  % avoid overfull lines from long identifiers and paths

\newcommand{\code}[1]{\texttt{#1}}
\newcommand{\Ang}{\mbox{\AA}}  % math- and text-safe Angstrom
\newcommand{\pearson}{Pearson~$R$}

\title{\textbf{What a Protein--Ligand Affinity Model Actually Learns:\\
A Preprocessing Bug, Three Null Ablations, and a\\
Proximity Ceiling on CASF-2016}}

\author{Pankaj Kumar\thanks{Independent researcher. Correspondence: \texttt{pankajator@gmail.com}.
This study was carried out with substantial assistance from an AI coding agent, which
wrote the implementation, ran the experiments, and drafted the manuscript under the
author's direction. All quantitative claims are reproducible from the committed code and
run artifacts.}}
\date{November 2026}

\begin{document}
\maketitle

\begin{abstract}
We report a controlled empirical study of BindFormer, a residue--atom co-attention model
that predicts protein--ligand binding affinity ($\mathrm{p}K$) from a complex with a known
crystal pose, evaluated for scoring power on CASF-2016. The study is organized around an
evaluation protocol built \emph{before} any architectural change: every comparison uses
three random seeds and a paired bootstrap over the 285 test complexes, so that an effect
the size of a typical architectural tweak (0.01--0.03 \pearson{}) can be told apart from
seed and test-set noise. The protocol repays its cost immediately. Our single-seed
baseline reaches $R=0.748\pm0.011$, but trivial predictors that use no binding information
reach $R=0.50$ (ligand heavy-atom count) and $R=0.51$ (pocket residue count), so roughly
two thirds of the raw correlation is explained by molecular size and pocket identity. A
diagnostic baseline then exposed a train/test preprocessing bug---hydrogen atoms were
silently retained in the CASF ligands but not the training ligands---whose correction
raised \pearson{} by about $0.08$ everywhere and reversed two conclusions that had
otherwise looked significant, including an apparent $+0.14$ benefit of ESM-2 protein
embeddings that corrected to $+0.05$ at $p=0.095$. Of four architectural or objective
interventions, three produced no significant change on CASF and one significantly hurt,
each with an accompanying diagnostic confirming the intervention was mechanistically
active rather than inert. A pose-perturbation analysis explains the pattern: rotating the
ligand in place costs only $0.02$--$0.03$ \pearson{}, while translating it out of the
pocket drives the score to zero. BindFormer scores \emph{composition and proximity}, not
interaction geometry, and refining how attention weights distances cannot help a model
that does not read fine geometry. We offer the protocol, the diagnostics, and the
preprocessing lesson as transferable, and release all code.
\end{abstract}

\section{Introduction}

Machine-learning scoring functions for protein--ligand binding affinity are usually
reported as a single correlation coefficient on a benchmark such as the CASF-2016 scoring
power set~\cite{casf2016}: Pearson~$R$, RMSE, and the residual standard deviation SD over
285 protein--ligand complexes. Published deep models reach roughly $0.80$--$0.85$
\pearson{}. A single number on 285 complexes, however, carries a 95\% confidence interval
of about $\pm0.06$, which is wider than the improvement most architectural changes claim.
Reporting one number per configuration therefore cannot, on its own, distinguish a real
advance from noise, and the literature's difficulty in reproducing incremental gains is at
least partly a measurement problem.

This paper is a case study in treating the measurement as the first-class object. We build
BindFormer, a residue--atom co-attention affinity model in the current idiom (ESM-2 pocket
embeddings, bidirectional distance-masked co-attention, an auxiliary contact head), and we
build an uncertainty harness---multiple seeds plus a paired bootstrap over complexes---%
before attempting any improvement. We then use that harness to interrogate the model rather
than only to rank it.

The study's findings are mostly negative, and deliberately so:

\begin{enumerate}
\item A correctly measured baseline reaches $R=0.748\pm0.011$ (three seeds), with a
three-seed prediction ensemble at $0.772$.
\item Trivial, no-binding predictors reach $R=0.50$--$0.53$ on CASF-2016, so most of the
raw correlation requires no model of binding at all. This floor is rarely reported and
materially changes how a headline number should be read.
\item A diagnostic baseline uncovered a hydrogen-handling bug that silently inflated the
CASF ligands relative to the training ligands. Correcting it was worth $+0.08$ \pearson{},
larger than any architectural change we tried, and it reversed two conclusions.
\item Three architectural or objective interventions (a radial-basis distance bias,
attention-pooled readout, and a decoy-pose discrimination objective) produced no
significant change on CASF; a sharpened attention pool significantly \emph{reduced} it. In
each case a mechanism diagnostic confirms the intervention was active.
\item Pose perturbation shows the model is almost insensitive to ligand orientation:
rotating the ligand within its pocket costs $0.02$--$0.03$ \pearson{}. The model reads
proximity and composition, not interaction geometry.
\end{enumerate}

We do not claim a new state of the art. We claim that a disciplined protocol, paired with
cheap mechanistic diagnostics, turns a sequence of ambiguous results into a coherent and
falsifiable account of what the model has and has not learned---and catches at least one
spurious ``significant'' result and one data bug that a single-number workflow would have
shipped.

\section{Model and data}

\subsection{Architecture}

BindFormer scores a given crystal pose. Pocket residues within $10\,\Ang$ of the ligand are
encoded with ESM-2 650M per-residue embeddings~\cite{esm2}; ligand heavy atoms receive
74-dimensional atom features together with a learned per-head bias derived from bond-graph
shortest-path distances, in the manner of a graph transformer~\cite{graphormer}.
Bidirectional residue--atom co-attention is restricted to residue--atom pairs within
$10\,\Ang$. The model has two heads: affinity regression and an auxiliary residue--atom
contact head trained against $4\,\Ang$ labels. The encoders use a custom biased transformer
built on scaled dot-product attention rather than \code{nn.TransformerEncoder}, because the
PyTorch fast path turned the per-head bias mask into \code{NaN} for batch sizes greater than
one.

\subsection{Data and splits}

We use PDBbind v2020~\cite{pdbbind} (refined and general sets, 19{,}443 complexes) for
training and validation, and the CASF-2016 core set (285 complexes) as a held-out test set
that never enters training or validation. Validation is a seeded random 10\% of the
non-core data. Featurization crops each protein to its $10\,\Ang$ pocket, computes contact
labels, and caches per-complex tensors; protein features are stored in half precision.
Five of the 19{,}443 ligands fail all parsing fallbacks (corrupt SDF records in the general
set, none in CASF) and are dropped, leaving splits of 17{,}238 train / 1{,}915 validation /
285 CASF, verified to be disjoint complex-by-complex.

All training was run on a single Apple M5 Max using the Metal Performance Shaders backend.
Preprocessing with ESM-2 took about 11 minutes; each training seed took 70--150 minutes
depending on machine contention. The full suite comprises 50 unit tests.

\section{Evaluation protocol}

\subsection{The measurement problem}

On $n=285$ complexes, the 95\% confidence interval on a single \pearson{} near $0.75$ is
approximately $\pm0.06$ (bootstrap). The seed-to-seed standard deviation of our baseline is
$0.011$--$0.022$. The architectural changes we test are each expected, from the literature,
to move \pearson{} by $0.01$--$0.03$. All three of these quantities are the same order of
magnitude, so a comparison of two configurations by one number each is uninformative by
construction.

\subsection{The harness}

We therefore adopt a fixed protocol for every comparison (\code{scripts/compare.py}):

\begin{itemize}
\item \textbf{Seeds.} Each configuration is trained at three seeds into separate
directories; we report the per-seed mean and standard deviation.
\item \textbf{Paired bootstrap.} Two configurations are compared by resampling the 285
\emph{complexes} (10{,}000 resamples) and recomputing the metric difference on the
seed-averaged predictions, both models scored on the same complexes. Pairing removes
test-set sampling noise; seeds remove initialization noise.
\item \textbf{Residual correlation.} Pairing tightens the interval only in proportion to how
correlated the two models' residuals are. The tool reports that correlation rather than
assuming it: it is $\approx0.95$ for near-identical variants (where pairing roughly halves
the interval) and $\approx0.73$ for ESM-2 versus one-hot (where pairing buys little). We
report both the paired interval and an unpaired Fisher-$z$ reference.
\end{itemize}

A supporting tool, \code{scripts/eval\_checkpoints.py}, re-scores a saved checkpoint on a
new feature cache without retraining; this is what made the preprocessing-bug correction
(Section~\ref{sec:bug}) a matter of re-featurizing and re-scoring rather than repeating
every run.

\section{Results}

\subsection{Baseline and the no-binding floor}
\label{sec:baseline}

Three ESM-2 seeds score CASF-2016 at $0.751$, $0.758$, $0.736$ \pearson{}
($0.748\pm0.011$); the three-seed prediction ensemble reaches $0.772$ with a bootstrap
interval $[0.705, 0.828]$ and RMSE $1.44$. Against the $0.80$--$0.85$ range of published
deep models this is a reasonable but unremarkable baseline.

The headline is better read against a floor. Table~\ref{tab:trivial} reports predictors
that use no binding information, fitted on the training split and scored on CASF.

\begin{table}[t]
\centering
\begin{tabular}{lcc}
\toprule
Predictor & CASF \pearson{} & RMSE \\
\midrule
Ligand heavy-atom count (no fit)     & 0.500 & 2.12 \\
Pocket residue count (no fit)        & 0.508 & 1.98 \\
Ridge on per-element counts          & 0.526 & 1.91 \\
Ridge on mean/sum atom features      & 0.526 & 1.91 \\
\textbf{BindFormer (per seed)}       & \textbf{0.748} & \textbf{1.47} \\
\bottomrule
\end{tabular}
\caption{No-binding baselines on CASF-2016. Simply counting ligand atoms reaches
$R=0.50$; counting pocket residues reaches $0.51$. BindFormer adds about $0.22$ over this
floor.}
\label{tab:trivial}
\end{table}

Molecular size alone reaches $R=0.50$, and pocket size $0.51$. Roughly two thirds of the
raw correlation on CASF-2016 requires no binding model. We report this not to diminish the
benchmark but because it reframes every subsequent number: a change that moves \pearson{}
by $0.02$ is moving a quantity whose first $0.50$ is molecular bookkeeping.

\subsection{A train/test preprocessing bug}
\label{sec:bug}

The trivial baselines were introduced as a diagnostic, and they immediately exposed a bug.
Ridge regression on per-element atom counts scored $0.04$---below heavy-atom count's
$0.46$---which is impossible, since element counts are a superset of the total count. The
cause was a single element column with mean $0.11$ in training and $13.96$ in CASF:
hydrogen.

The CASF-2016 core-set SDF files carry explicit hydrogens; the PDBbind general and refined
SDFs do not. RDKit's \code{removeHs=True} preserves hydrogens it deems chemically
meaningful rather than stripping them unconditionally, so $36\%$ of cached CASF ligand
atoms were hydrogen while the training ligands were heavy-atom only. The test ligands were
about $1.6\times$ the length of the training ligands (median $36$ versus $23$ atoms), and
because the pocket crop and the $4\,\Ang$ contact labels are derived from ligand
coordinates, those were computed against hydrogen positions as well. The model was trained
on one atom distribution and evaluated on another.

The fix is one call to \code{RemoveAllHs} on every ligand parse path---the behavior the
featurizer's documentation already claimed. Because the training data was almost unaffected
($\approx11\%$ of complexes lost $1$--$3$ stray hydrogens, verified to preserve finiteness
and atom counts), the correction required only re-featurizing the test set and re-scoring
the existing checkpoints. It raised CASF \pearson{} by about $0.08$ for every
configuration and reversed two conclusions (Sections~\ref{sec:esm} and~\ref{sec:ablations}).
It was the single largest improvement in the project, larger than any architectural change
attempted, which is the study's central cautionary result.

\subsection{Protein features: ESM-2 versus one-hot}
\label{sec:esm}

Replacing ESM-2 pocket embeddings with a one-hot residue-identity encoding, on identical
splits, gives the corrected comparison in Table~\ref{tab:esm}.

\begin{table}[t]
\centering
\begin{tabular}{lcc}
\toprule
Protein features & CASF \pearson{} & val \pearson{} \\
\midrule
ESM-2 650M (3 seeds) & 0.748 & 0.741--0.751 \\
one-hot (1 seed)     & 0.722 & 0.727 \\
\bottomrule
\end{tabular}
\caption{Corrected ESM-2 versus one-hot. Paired bootstrap on seed-mean predictions:
$\Delta R = -0.050\ [-0.108, +0.009]$, $p=0.095$---not significant.}
\label{tab:esm}
\end{table}

Measured through the hydrogen bug, this comparison read as $\Delta R = -0.154$ at
$p=0.0002$: an apparently decisive $+0.14$ \pearson{} benefit from ESM-2. The corrupted,
oversized CASF ligands harmed the weaker one-hot model (which fell to $0.551$) far more than
the ESM-2 model (which fell to $0.670$), inflating the gap roughly threefold. Corrected, the
gap is $-0.050$ at $p=0.095$: ESM-2 may be worth about $0.05$ \pearson{} here, but on three
seeds it is \emph{not established}. We withdraw the original claim. A related claim---that
validation compressed the ESM-2 gap ``tenfold''---is also withdrawn; corrected, the gap is
$0.019$ on validation and $0.026$ on CASF, so validation was a fair proxy for this
comparison.

One validation finding survives the correction: validation does not rank \emph{seeds}.
Across the three baseline seeds, the best validation \pearson{} belongs to the worst CASF
seed and vice versa; validation \pearson{} moves $0.010$ across seeds while CASF moves
$0.022$. Validation remains usable for early stopping but not for choosing among
configurations or checkpoints.

\subsection{Architectural ablations}
\label{sec:ablations}

We tested two architectural changes that target the model's use of geometry, each on three
seeds against the three baseline seeds, each paired on the 285 complexes. For each we
computed a diagnostic establishing the mechanism was active, so that a null result means
``did not help'' rather than ``did nothing.''

\paragraph{Radial-basis distance bias.} Pairwise distances (residue--residue in protein
self-attention; residue--atom in co-attention) are expanded on 16 Gaussian radial basis
functions and mixed into one additive logit per attention head. The output projection is
zero-initialized, so an untrained model is numerically identical to the baseline (asserted
to $10^{-6}$), and any effect is learned. Result: $\Delta R = -0.003\ [-0.024, +0.018]$,
$p=0.77$---no effect. The mechanism was active: sweeping distance through the trained
modules yields a co-attention bias spanning $1.4$--$2.6$ attention logits, enough to reshape
the softmax substantially.

\paragraph{Attention-pooled readout.} The affinity head pools residues and atoms with a
masked mean; we replaced this with a learned attention pool, again zero-initialized so that
uniform weights reproduce the mean exactly. A first attempt produced $\Delta R = -0.004$ at
$p=0.70$, but a diagnostic showed the pooling had never left uniform (attention entropy
$0.999$ of uniform), making the null uninformative. We then forced the pool to train (an
MLP scorer, softmax temperature $0.5$, and a $20\times$ learning rate for the scorer),
verified on a short run that attention entropy fell to $0.67$--$0.81$, and re-ran. The
sharpened pool reaches attention entropy $0.05$--$0.51$ of uniform at the selected
checkpoints (the top-weighted token taking $23$--$58\times$ its uniform share) and scores
$\Delta R = -0.032\ [-0.056, -0.010]$, $p=0.0076$: a significant \emph{regression}, with
RMSE agreeing ($+0.076$, $p=0.004$).

That mean pooling and near-argmax pooling score within $0.004$ of each other---and that
forcing selectivity only hurts---indicates the readout is not where the limiting information
lives. The optimizer traverses the whole path from uniform to near-argmax on its own and
performance is flat along it.

\subsection{What the model uses: pose sensitivity}
\label{sec:pose}

If the pocket crop and co-attention mask are built from the true pose, a model can score
well while ignoring geometry. We test this directly by perturbing the ligand pose at
evaluation time, holding all protein and ligand \emph{tokens} fixed so that only the
geometry entering the distance mask (and the RBF bias, where present) changes
(\code{scripts/pose\_sensitivity.py}). Table~\ref{tab:pose} reports the sweep, averaged over
three seeds.

\begin{table}[t]
\centering
\begin{tabular}{lccc}
\toprule
Pose & \pearson{} & RMSE & mean $|\Delta\mathrm{pred}|$ (p$K$) \\
\midrule
true                 & 0.748 & 1.49 & --- \\
rotated in place     & 0.728 & 1.56 & 0.27 \\
shifted $2\,\Ang$    & 0.743 & 1.53 & 0.28 \\
shifted $5\,\Ang$    & 0.677 & 1.81 & 0.66 \\
shifted $10\,\Ang$   & 0.324 & 3.20 & 1.63 \\
shifted $25\,\Ang$   & $-0.096$ & 17.77 & 12.27 \\
scrambled            & 0.309 & 3.25 & 1.60 \\
\bottomrule
\end{tabular}
\caption{Pose sensitivity on CASF-2016. ``Rotated in place'' keeps the ligand in its
pocket, keeps its centroid, and keeps its internal shape; only which atom faces which
residue changes. It costs about $0.02$ \pearson{}. Translating the ligand out of the pocket
destroys the score, but drives the model out of distribution (at $25\,\Ang$ the predictions
are nonsense, not merely wrong), so those rows show the model needs the ligand in the
pocket, not that it models interactions.}
\label{tab:pose}
\end{table}

The informative row is the rotation. Randomizing which atom contacts which residue, while
leaving the pocket, the centroid, and the ligand shape untouched, costs only $0.02$
\pearson{}; a paired bootstrap against the true pose gives $-0.018\ [-0.032, -0.004]$
($p=0.011$) for a baseline trained on the pre-correction data and $-0.032\ [-0.050, -0.016]$
for a baseline trained on corrected data. Either way, atom-level pairing geometry carries
almost none of the model's scoring power. The large drops under translation conflate
``geometry matters'' with ``the input left the training distribution,'' and should not be
read as evidence that interactions are modeled.

Together with the trivial floor (Section~\ref{sec:baseline}), this gives a coherent account:
BindFormer scores composition and proximity---the right sort of ligand in the right sort of
pocket---and barely uses interaction geometry. This also explains the ablations of
Section~\ref{sec:ablations}: refining how attention weights distances or pools tokens cannot
help a model that is not reading fine geometry in the first place.

\subsection{A pose-discrimination objective}
\label{sec:decoy}

If affinity regression gives the encoder no reason to use geometry, perhaps an explicit
objective will. We added a native-versus-decoy discrimination head whose signal is
geometry alone: a decoy shares the native pose's protein, atom features, and bond graph,
differing only in the ligand's placement (a random rotation in place, exactly the
perturbation the baseline was insensitive to). The decoy's cached attention mask is dropped
so the model must rebuild it from the moved coordinates, and a short run confirmed the task
is learnable (native-versus-decoy ranking accuracy rising from chance to $76\%$ in three
epochs) before committing seeds.

We trained three seeds with the objective and, as a matched control, three baseline seeds on
the identical corrected cache---so that the comparison is not confounded by the preprocessing
fix. Results in Table~\ref{tab:decoy}.

\begin{table}[t]
\centering
\begin{tabular}{lccc}
\toprule
Configuration & per-seed CASF \pearson{} & seed SD & ensemble \\
\midrule
+ decoy objective                 & 0.756 & 0.022 & 0.779 \\
baseline, corrected training       & 0.733 & 0.017 & 0.767 \\
baseline, pre-fix training (rescored) & 0.748 & 0.011 & 0.772 \\
\bottomrule
\end{tabular}
\caption{Decoy-pose objective. Paired bootstrap on seed-mean predictions: versus the
matched control, $\Delta R = +0.012\ [-0.008, +0.033]$, $p=0.23$; versus the rescored
baseline, $+0.007$, $p=0.48$. Not significant against either.}
\label{tab:decoy}
\end{table}

The objective is not significant against either control. Its point estimate is positive
(unlike the RBF and pooling ablations) and the upper end of its interval ($+0.033$) would be
worth having, so it is the only intervention that might survive more seeds; it is not
established on three. More telling is the mechanism check. The pose head learns its task
well---ranking accuracy $83.6\%$, $84.6\%$, $88.3\%$ at the selected epochs---yet the paired
rotation penalty of the affinity head is unchanged: $-0.030\ [-0.052, -0.009]$ with the
objective versus $-0.032\ [-0.050, -0.016]$ without it. The seed with the \emph{best} pose
ranking ($88.3\%$) had the \emph{worst} CASF score ($0.731$). Learning to discriminate poses
better did not make the affinity head more pose-aware; the auxiliary task and the regression
head are decoupled.

This result also forced a mid-study correction. An earlier reading had compared the
pose-trained models' rotation penalty against baselines trained on the pre-correction data
($-0.018$) and reported the objective as roughly doubling pose sensitivity. The matched
baseline on corrected data already sits at $-0.032$, so the increase is attributable to the
hydrogen fix---training on hydrogen-free ligands makes the model read placement more
sharply---and not to the decoy objective. We note it as a second instance of the matched
control, rather than a convenient one, changing the conclusion.

\section{Discussion}

\subsection{A scorecard of interventions}

Table~\ref{tab:scorecard} collects the interventions. The one data-preprocessing fix
outweighs the four architectural and objective changes combined; three of the four produced
no significant effect and one was harmful. The ``mechanism active'' column is what makes the
nulls informative: in every case we can show the intervention did what it claimed, so the
null is a statement about the model, not about a dead code path.

\begin{table}[t]
\centering
\begin{tabular}{lccl}
\toprule
Intervention & $\Delta R$ (CASF) & $p$ & mechanism verified active \\
\midrule
ESM-2 vs.\ one-hot            & $+0.050$ & 0.095 & --- \\
RBF distance bias             & $-0.003$ & 0.77  & yes ($1.4$--$2.6$ logit span) \\
Attention pool (as built)     & $-0.004$ & 0.70  & no (entropy $0.999$) \\
Attention pool (sharpened)    & $-0.032$ & 0.008 & yes (entropy $0.05$--$0.51$) \\
Decoy-pose objective          & $+0.012$ & 0.23  & task learned; head unchanged \\
\textit{Hydrogen preprocessing fix} & \textit{$+0.08$} & --- & --- \\
\bottomrule
\end{tabular}
\caption{All interventions, measured on CASF-2016 by paired bootstrap over 285 complexes,
three seeds per arm. The preprocessing fix is not an architectural change and has no paired
$p$; it is the correction underlying every other row's corrected value.}
\label{tab:scorecard}
\end{table}

\subsection{Spurious results caught by the protocol}

Three conclusions that a single-number workflow would have reported did not survive the
protocol: the $+0.14$ ESM-2 benefit (an artifact of the preprocessing bug), the claim that
validation-based early stopping selects a better CASF checkpoint (within noise, $p=0.17$),
and the claim that the decoy objective doubled pose sensitivity (attributable to the data
fix via the matched control). Each was caught by a control or a mechanism diagnostic, not
by re-examination of the headline number. This is the practical argument for the protocol:
at the noise scale of CASF-2016, cheap controls are not optional rigor but the difference
between a result and an artifact.

\subsection{Implications for architecture search}

Our results suggest that, for a model in this family, searching over attention-level
inductive biases is poorly targeted. The model's ceiling is not set by how it weights
distances inside attention but by what its training objective rewards: affinity regression
on crystal poses rewards reading the pocket and the ligand, and penalizes nothing about
getting the geometry wrong. The decoy experiment shows the model \emph{can} be trained to
read pose, but as a side task that does not reach the regression head. The productive next
step is to couple the geometric signal to the affinity prediction---sharing the
pose-discrimination representation with the readout, or weighting the affinity loss by pose
confidence---rather than adding attention capacity.

\section{Limitations}

Several limitations bound these claims. The one-hot ablation uses a single seed, so its
comparison, while corrected, has no error bar on the weaker arm. The decoy objective's two
controls differ by $0.015$ \pearson{} per seed (most plausibly noise between independent
runs, since the preprocessing fix was verified not to damage the training data), which
leaves its effect estimate sensitive to the choice of control; this is itself a reason to
regard it as unresolved rather than as evidence of no effect. All evaluation uses the
standard random CASF split; we have not run the harder protein-family or LP-PDBBind splits,
so we cannot say how much of $0.748$ is train/test similarity. The model scores a single
crystal pose, and the auxiliary contact head, though trained, is never evaluated against
reference interaction fingerprints. Finally, the absolute baseline ($0.748$) is below
published state of the art; the study is about measurement and mechanism, not about topping
the benchmark.

\section{Conclusion}

We built a protein--ligand affinity model and, more importantly, an evaluation protocol to
go with it, and used the protocol to ask what the model learns rather than only how high it
scores. The protocol paid for itself: it exposed a preprocessing bug worth more than any
architectural change we attempted, caught two further spurious conclusions, and turned three
null ablations into a single coherent finding---that the model scores composition and
proximity, not interaction geometry. We recommend that scoring-power results on CASF-2016 be
reported against the no-binding floor ($R\approx0.50$), with multiple seeds and paired
comparisons, and with a mechanism diagnostic for any claimed architectural benefit. The
cost is a few extra training runs; the return is the ability to tell a result from an
artifact at a noise scale where the two look identical.

\paragraph{Reproducibility.} All code, the uncertainty harness, the diagnostic scripts, and
the per-seed run records are available at
\url{https://github.com/pankajlp/bindformer-hierarchical-residue-atom-co-attention}. Every
quantitative claim in this paper is reproducible from the committed code; \code{RESULTS.md}
in the repository records the raw per-seed numbers and the exact commands. Trained model weights are not redistributed with the code but are available from the author on request; their sha256 sums are committed so any copy can be verified.

\begin{thebibliography}{9}

\bibitem{casf2016}
M.~Su, Q.~Yang, Y.~Du, G.~Feng, Z.~Liu, Y.~Li, and R.~Wang.
\newblock Comparative assessment of scoring functions: the CASF-2016 update.
\newblock \emph{Journal of Chemical Information and Modeling}, 59(2):895--913, 2019.

\bibitem{pdbbind}
Z.~Liu, M.~Su, L.~Han, J.~Liu, Q.~Yang, Y.~Li, and R.~Wang.
\newblock Forging the basis for developing protein--ligand interaction scoring functions.
\newblock \emph{Accounts of Chemical Research}, 50(2):302--309, 2017.

\bibitem{esm2}
Z.~Lin, H.~Akin, R.~Rao, B.~Hie, Z.~Zhu, W.~Lu, et al.
\newblock Evolutionary-scale prediction of atomic-level protein structure with a language
model.
\newblock \emph{Science}, 379(6637):1123--1130, 2023.

\bibitem{graphormer}
C.~Ying, T.~Cai, S.~Luo, S.~Zheng, G.~Ke, D.~He, Y.~Shen, and T.-Y. Liu.
\newblock Do transformers really perform badly for graph representation?
\newblock In \emph{Advances in Neural Information Processing Systems (NeurIPS)}, 2021.

\end{thebibliography}

\end{document}
